% equations.tex — auto-generated by extract_equations.py; do not edit.
%
% Every display equation of "Investing in Artificial General Intelligence"
% (main text and Internet Appendix), with its introducing prose and its
% Quarto label (shown as [eq-...] next to the equation). The labels match
% those referenced in the Lean 4 formalization's module docstrings and in
% lean/README.md, so this file is the reference against which the Lean
% package can be checked for faithfulness.
\documentclass[11pt]{article}
\usepackage{amsmath,amssymb}
\usepackage[margin=1in]{geometry}
\usepackage[hidelinks]{hyperref}
\setlength{\parindent}{0pt}
\setlength{\parskip}{6pt}
\title{Investing in Artificial General Intelligence:\\ complete equation listing}
\author{Auto-generated from the paper's Quarto sources}
\date{}
\begin{document}
\maketitle
\noindent\emph{Each display equation from the paper and the Internet Appendix
appears below in source order, preceded by the sentence(s) that introduce it
and followed by the immediate explanation of its symbols. Bracketed monospace
tags (e.g.\ {[\texttt{eq-trigger-phi}]}) are the paper's equation labels; the
Lean formalization references these labels. Cross-references and citation keys
from the manuscript are kept as monospace tokens.}
\bigskip

\section*{The Model}

\subsection*{Environment}

Time is continuous. A stochastic demand process $X_t$ governs the revenue available to frontier AI labs. Demand follows a geometric Brownian motion (GBM) with regime-dependent drift and common volatility:

\begin{equation*}
\tag*{[\texttt{eq-gbm}]}
dX_t = \mu_s X_t \, dt + \sigma X_t \, dW_t, \quad s \in \{L, H\},
\end{equation*}

where $W_t$ is a standard Brownian motion, $s$ indexes the demand regime, and $\sigma > 0$ is the common volatility parameter.



\subsection*{Technology}

\subsubsection*{Training-inference allocation}

- \textbf{Low regime} ($L$, pre-AGI): revenue derives from \emph{inference}---serving current demand for AI predictions and services. Firm $i$'s flow revenue is:

\begin{equation*}
\tag*{[\texttt{eq-revenue-L}]}
\pi_i^L(X) = X \cdot [(1-\phi_i)K_i]^\alpha,
\end{equation*}

where $\alpha \in (0,1)$ captures diminishing returns to inference capacity.



- \textbf{High regime} ($H$, post-AGI): revenue derives from \emph{training quality}---the capability of the firm's models determines its competitive position in a world of sustained high demand. Firm $i$'s flow revenue is:

\begin{equation*}
\tag*{[\texttt{eq-revenue-H}]}
\pi_i^H(X) = X \cdot (\phi_i K_i)^\alpha.
\end{equation*}



\subsection*{Single-Firm Benchmark}

\subsubsection*{Installed value}

In the absorbing high regime, the value of installed capacity depends on training quality:

\begin{equation*}
\tag*{[\texttt{eq-installed-value-H}]}
V_H(X, K, \phi) = \frac{X (\phi K)^\alpha}{r - \mu_H} - \frac{\delta K}{r} = A_H \cdot X \cdot (\phi K)^\alpha - \frac{\delta K}{r},
\end{equation*}

where $A_H = 1/(r - \mu_H)$. Revenue depends on training compute $\phi K$: a firm that allocated all capacity to inference ($\phi = 0$) generates no H-regime revenue, while a firm that invested heavily in training captures more value in the high-demand world.



In the low regime, the value of installed capacity combines \emph{current inference revenue} and the \emph{expected continuation value} from a potential regime switch to $H$:

\begin{equation*}
\tag*{[\texttt{eq-installed-value-L}]}
V_L(X, K, \phi) = A_{\text{eff}}(\phi, K) \cdot X - \frac{\delta K}{r},
\end{equation*}



where the \emph{effective revenue coefficient} is:

\begin{equation*}
\tag*{[\texttt{eq-a-eff}]}
A_{\text{eff}}(\phi, K) = \frac{[(1-\phi)K]^\alpha}{r - \mu_L + \lambda} + \frac{\lambda}{r - \mu_L + \lambda} \cdot \frac{(\phi K)^\alpha}{r - \mu_H}.
\end{equation*}



\subsubsection*{Option value in regime $H$}

In the absorbing high regime, the option to invest at trigger $X_H^*$ with capacity $K_H^*$ and training fraction $\phi_H^*$ has the standard real options form (\texttt{dixit1994investment}):

\begin{equation*}
\tag*{[\texttt{eq-option-H}]}
F_H(X) = \left(V_H(X_H^*, K_H^*, \phi_H^*) - I(K_H^*)\right) \left(\frac{X}{X_H^*}\right)^{\beta_H}, \quad X < X_H^*,
\end{equation*}



where $\beta_H > 1$ is the positive root of the H-regime characteristic equation:

\begin{equation*}
\tag*{[\texttt{eq-beta-H}]}
\tfrac{1}{2}\sigma^2 \beta(\beta - 1) + \mu_H \beta - r = 0.
\end{equation*}



The optimal trigger satisfies the smooth-pasting condition:

\begin{equation*}
\tag*{[\texttt{eq-trigger-H}]}
X_H^* = \frac{\beta_H}{\beta_H - 1} \cdot \frac{(r - \mu_H)}{(\phi K)^{\alpha}} \left(c K^\gamma + \frac{\delta K}{r}\right),
\end{equation*}

which shows the standard result that the trigger exceeds the Marshallian (NPV = 0) level by the factor $\beta_H / (\beta_H - 1) > 1$---the \emph{option premium} that reflects the value of waiting.



\textbf{Proposition 1} (Optimal capacity and training fraction; single-firm benchmark). *The firm jointly chooses capacity $K^*$ and training fraction $\phi^*$ to maximize the option value. The optimal capacity is:*

\begin{equation*}
K^* = \left[\frac{\delta\,(\alpha \beta_H - \beta_H + 1)}{r\,c\,(\gamma(\beta_H - 1) - \alpha \beta_H)}\right]^{\frac{1}{\gamma - 1}},
\end{equation*}

*provided $1/\gamma < (\beta_H - 1)/(\alpha \beta_H) < 1$. Note that $K^*$ is independent of $\phi$: the effective revenue coefficient takes the form $A_{\text{eff}} = g(\phi) \cdot K^\alpha$ (both in the pure H-regime and under regime switching), so the $\phi$-dependent factor cancels in the first-order condition for $K$ (see Internet Appendix A, Step 4). The optimal $\phi^*$ is then determined separately by maximizing $A_{\text{eff}}$ over $\phi$, and the joint optimum $(K^*, \phi^*)$ yields the investment trigger via smooth pasting.*



\subsubsection*{Option value in regime $L$}

In the low regime, the firm's investment option reflects both the possibility of investing in regime $L$ and the potential for a regime switch to $H$. The option value satisfies the Hamilton-Jacobi-Bellman equation incorporating the regime-switching term:

\begin{equation*}
\tag*{[\texttt{eq-hjb-L}]}
\tfrac{1}{2}\sigma^2 X^2 F_L'' + \mu_L X F_L' + \lambda\bigl[F_H(X) - F_L(X)\bigr] - r\, F_L = 0,
\end{equation*}



where $F_H(X) = B_H X^{\beta_H}$ is the known H-regime \emph{option} value (\texttt{eq-option-H})---the value of the investment option conditional on being in regime $H$ without having invested yet, not the installed value $V_H$. Substituting and rearranging, this becomes a second-order Euler ODE with non-homogeneous term $\lambda B_H X^{\beta_H}$:

\begin{equation*}
\tfrac{1}{2}\sigma^2 X^2 F_L'' + \mu_L X F_L' - (r + \lambda)\, F_L + \lambda B_H X^{\beta_H} = 0.
\end{equation*}



The general solution has two components. The \emph{homogeneous} solution has characteristic equation $\tfrac{1}{2}\sigma^2 \beta(\beta - 1) + \mu_L \beta - (r + \lambda) = 0$, with positive root $\beta_L^+$ that depends on $(\sigma, \mu_L)$ and the effective discount rate $r + \lambda$. The \emph{particular} solution, of the form $C \cdot X^{\beta_H}$, arises from the regime-switching term, with coefficient:

\begin{equation*}
\tag*{[\texttt{eq-particular-C}]}
C = \frac{-\lambda\, B_H}{Q_L(\beta_H)}, \quad Q_L(\beta) \equiv \tfrac{1}{2}\sigma^2 \beta(\beta - 1) + \mu_L \beta - (r + \lambda).
\end{equation*}



Applying the boundary condition $F_L(0) = 0$ (which eliminates the negative root) yields:

\begin{equation*}
\tag*{[\texttt{eq-option-L-general}]}
F_L(X) = A_1 X^{\beta_L^+} + C \cdot X^{\beta_H}, \quad X < X^*.
\end{equation*}



Under this convention, the investment trigger for the joint $(K^*, \phi^*)$ optimization satisfies:

\begin{equation*}
\tag*{[\texttt{eq-trigger-phi}]}
X^* = \frac{\beta_H}{\beta_H - 1} \cdot \frac{\delta K^*/r + c(K^*)^\gamma}{A_{\text{eff}}(\phi^*, K^*)},
\end{equation*}

which generalizes \texttt{eq-trigger-H} by replacing the H-regime revenue coefficient with the effective coefficient $A_{\text{eff}}$ that combines both regimes. Note that \texttt{eq-trigger-phi} is an implicit equation: $A_{\text{eff}}(\phi^*, K^*)$ itself depends on the optimal policy, so the trigger is determined jointly with $(K^*, \phi^*)$ through the system of first-order conditions (Internet Appendix A). The option premium $\beta_H / (\beta_H - 1) > 1$ still applies, reflecting the value of waiting, but the relevant revenue metric now accounts for both L-regime inference earnings and the expected H-regime training prize.



\subsection*{Duopoly with Default Risk}

\subsubsection*{Regime-specific competition}

Revenue in each regime is governed by a Tullock contest success function (\texttt{tullock1980efficient; \texttt{skaperdas1996contest}}) over the regime-relevant capacity measure. The Tullock specification captures the essential feature of AI compute competition: firms' revenues depend on \emph{relative} capacity, not just absolute capacity, because customers (enterprises, developers) allocate workloads to providers based on relative performance and availability. In the L-regime, firms compete over \emph{inference capacity}:

\begin{equation*}
\tag*{[\texttt{eq-contest-L}]}
\pi_i^L(X, \mathbf{K}, \boldsymbol{\phi}) = X \cdot \frac{[(1-\phi_i)K_i]^{2\alpha}}{[(1-\phi_i)K_i]^\alpha + [(1-\phi_j)K_j]^\alpha},
\end{equation*}



which I write equivalently as $\pi_i^L = X \cdot [(1-\phi_i)K_i]^\alpha \cdot s_i^L$, where $s_i^L = [(1-\phi_i)K_i]^\alpha / \{[(1-\phi_i)K_i]^\alpha + [(1-\phi_j)K_j]^\alpha\}$ is the L-regime contest share. In the H-regime, firms compete over \emph{training quality}:

\begin{equation*}
\tag*{[\texttt{eq-contest-H}]}
\pi_i^H(X, \mathbf{K}, \boldsymbol{\phi}) = X \cdot \frac{(\phi_i K_i)^{2\alpha}}{(\phi_i K_i)^\alpha + (\phi_j K_j)^\alpha},
\end{equation*}

with H-regime contest share $s_i^H = (\phi_i K_i)^\alpha / \{(\phi_i K_i)^\alpha + (\phi_j K_j)^\alpha\}$.



\subsubsection*{Installed value in the duopoly}

The installed value for firm $i$ in the L-regime combines inference revenue and the expected H-regime continuation:

\begin{equation*}
\tag*{[\texttt{eq-duopoly-value-L}]}
V_i^L(X) = A_{\text{eff},i} \cdot X - \frac{\delta K_i}{r},
\end{equation*}



where the effective revenue coefficient for firm $i$ is:

\begin{equation*}
\tag*{[\texttt{eq-a-eff-duopoly}]}
A_{\text{eff},i} = \frac{[(1-\phi_i)K_i]^\alpha \cdot s_i^L}{r - \mu_L + \lambda} + \frac{\lambda}{r - \mu_L + \lambda} \cdot \frac{(\phi_i K_i)^\alpha \cdot s_i^H}{r - \mu_H}.
\end{equation*}



In the absorbing H-regime:

\begin{equation*}
V_i^H(X) = A_H \cdot X \cdot (\phi_i K_i)^\alpha \cdot s_i^H - \frac{\delta K_i}{r}.
\end{equation*}



\subsubsection*{Capital structure}

\subsubsection*{Default boundary}

Equity holders default endogenously when it is optimal to stop servicing debt. Following \texttt{leland1994corporate} and \texttt{leland1996optimal}, the optimal default boundary is obtained by smooth pasting. For a levered firm with capacity $K_i$, training fraction $\phi_i$, facing a rival with $(K_j, \phi_j)$:

\begin{equation*}
\tag*{[\texttt{eq-default-boundary}]}
X_D = \frac{\beta_s^-}{\beta_s^- - 1} \cdot \frac{c_D / r + \delta K_i / r}{A_{\text{eff},i}},
\end{equation*}

where $\beta_s^-$ (with $s = L$) is the negative root of the L-regime characteristic equation $\tfrac{1}{2}\sigma^2 \beta(\beta - 1) + \mu_L \beta - (r + \lambda) = 0$, incorporating the regime-switching term through the effective discount rate $r + \lambda$; $c_D$ is the coupon payment; and $A_{\text{eff},i}$ is the effective revenue coefficient from \texttt{eq-a-eff-duopoly.} The factor $\beta_s^- / (\beta_s^- - 1) \in (0, 1)$ reflects the option value of waiting before defaulting, yielding a lower (optimal) default boundary than the naive break-even condition.



*(ii) $X_D$ is decreasing in $\lambda$ precisely when the training allocation is sufficiently high, and the characterization is fully closed-form in the symmetric-shares case ($s_i^L = s_i^H$, as under symmetric capacities or monopoly). Two channels operate. The $A_{\text{eff}}$-channel (higher $\lambda$ raises the weight on the H-regime continuation value, lowering $X_D$) is positive if and only if $\phi_i$ exceeds*

\begin{equation*}
\tag*{[\texttt{eq-phi-underbar}]}
\underline{\phi} = \frac{R}{1 + R}, \quad R \equiv \left(\frac{r - \mu_H}{r - \mu_L}\right)^{1/\alpha}.
\end{equation*}



*An opposing markup channel operates through the discount root: higher $\lambda$ raises the effective discount rate and pushes $\beta_s^-/(\beta_s^- - 1)$ toward one, raising $X_D$. Because the $A_{\text{eff}}$-channel semi-elasticity is linear in the L-regime flow and H-regime continuation terms, the net condition $\partial X_D / \partial \lambda < 0$ also solves in closed form: it holds if and only if $\phi_i$ exceeds the exact threshold*

\begin{equation*}
\tag*{[\texttt{eq-phi-tilde}]}
\tilde{\phi} = \frac{\tilde{R}}{1 + \tilde{R}}, \quad \tilde{R} \equiv \left[\frac{(r - \mu_H)\,(1 + m\Delta)}{(r - \mu_L) - m\lambda\Delta}\right]^{1/\alpha},
\end{equation*}

*where $\Delta = r - \mu_L + \lambda$ and $m = \mathrm{d}\ln M / \mathrm{d}\lambda = 1/[\beta_s^-(\beta_s^- - 1)D] > 0$ is the markup-channel semi-elasticity, with $D = \sqrt{(\mu_L - \sigma^2/2)^2 + 2\sigma^2(r+\lambda)}$ (Internet Appendix A states the mild regularity condition). The exact threshold nests $\underline{\phi}$ as the $m \to 0$ limit and satisfies $\tilde{\phi} > \underline{\phi}$. For the baseline calibration ($r = 0.12$, $\mu_H = 0.06$, $\mu_L = 0.01$, $\alpha = 0.40$), $R = (0.06/0.11)^{2.5} \approx 0.22$, $\underline{\phi} \approx 0.18$, and $\tilde{\phi} \approx 0.32$; the optimal allocation $\phi^* \approx 0.70$ exceeds both, so faith-based survival holds at the optimum. Unlike $\underline{\phi}$, $\tilde{\phi}$ depends on $\lambda$: at the baseline calibration $\phi^*(\lambda) > \tilde{\phi}(\lambda)$ if and only if $\lambda \gtrsim 0.034$, so for sufficiently pessimistic beliefs the optimal allocation is too inference-heavy for the continuation-value channel to dominate, and $X_D$ is locally increasing in $\lambda$ at the optimum. Note that $\underline{\phi}$ is increasing in $\alpha$ (with less concavity in the revenue function, a larger training allocation is needed for the H-regime term to dominate), decreasing in $\mu_H - \mu_L$ (a larger regime-switch growth premium amplifies the H-regime term, lowering the threshold), and increasing in $r$ (higher discounting reduces the present value of the continuation opportunity).*



\subsubsection*{Equity and debt values}

Above the default boundary:

\begin{equation*}
E(X) = A_{\text{eff},i}\, X - \frac{\delta K}{r} - (1 - \ell) I(K) - \frac{c_D}{r} + \left[\frac{c_D}{r} + \frac{\delta K}{r} - A_{\text{eff},i}\, X_D \right]\left(\frac{X}{X_D}\right)^{\beta_s^-},
\end{equation*}

where $A_{\text{eff},i}$ is the effective revenue coefficient from \texttt{eq-a-eff-duopoly} and $\beta_s^- < 0$ is the negative characteristic root. The last term is the default option: it captures the value of limited liability, discounted by the probability-weighted factor $(X/X_D)^{\beta_s^-}$ which approaches one as $X \to X_D$. The bracket $[c_D/r + \delta K/r - A_{\text{eff},i}\, X_D]$ equals the ongoing flow deficit at the default boundary; the smooth-pasting conditions on the going-concern equity $E(X) + (1-\ell)I(K)$ pin the default boundary $X_D$ in \texttt{eq-default-boundary} and are independent of the sunk equity contribution. At $X = X_D$, $E(X_D) = -(1-\ell)I(K) \leq 0$: the equity holders' sunk contribution is lost, but limited liability prevents negative payoffs, so equity is worth $\max\{E(X),\, 0\}$. Without debt ($\ell = 0$), the expression simplifies to the NPV of the unlevered investment: $E(X) = A_{\text{eff},i}\, X - \delta K / r - I(K)$, treated as a perpetuity with no endogenous abandonment option.



Debt value is:

\begin{equation*}
D(X) = \frac{c_D}{r}\left[1 - \left(\frac{X}{X_D}\right)^{\beta_s^-}\right] + R(X_D) \left(\frac{X}{X_D}\right)^{\beta_s^-},
\end{equation*}



where the recovery in default is

\begin{equation*}
R(X_D) = \min\left\{(1 - b)\,\Lambda(X_D),\; \frac{c_D}{r}\right\}, \quad \Lambda(X_D) = \frac{[(1-\phi_i)K_i]^\alpha \, s_i^L}{r - \mu_L + \lambda}\, X_D,
\end{equation*}

$b \in [0,1]$ is the fraction of liquidation value lost in bankruptcy, and $\Lambda(X_D)$ is the liquidation value of the firm's productive assets at the default boundary (gross of maintenance costs, following \texttt{leland1994corporate}). Two features of the recovery specification deserve comment. First, only the \emph{inference} business survives liquidation: the buyer acquires the capacity serving current demand and the associated L-regime revenue stream, capitalized at the effective rate $r - \mu_L + \lambda$. The H-regime continuation value---the faith-based component of $A_{\text{eff},i}$---is destroyed in bankruptcy, because the speculative post-AGI payoff requires the going concern (its researchers, models, and training program) and does not transfer to creditors in a liquidation. This asymmetry is central to the credit-risk analysis in \texttt{sec-valuation}: faith-based survival lowers the \emph{probability} of default, but faith is worthless to creditors \emph{in} default, so training allocation raises the loss given default. Second, the cap at $c_D/r$ imposes absolute priority: creditors cannot recover more in default than the default-free value of their coupon claim. Without the cap, a lightly levered firm---whose default boundary is driven by the operating cost $\delta K$ rather than the coupon---would receive a windfall at default, since the liquidation value of its assets can exceed its small debt claim.



\section*{Quantitative Implications}

\subsection*{Value Decomposition}

1. \textbf{Assets-in-place}: the value of existing installed capacity at current demand, incorporating both L-regime inference revenue and H-regime continuation value through $A_{\text{eff}}$.

\begin{equation*}
V_{\text{AIP}} = A_{\text{eff}}(\phi, K_{\text{installed}}) \cdot X - \frac{\delta K_{\text{installed}}}{r}.
\end{equation*}



2. \textbf{Capacity gap value}: the shortfall of the current installed base relative to the optimally sized project,

\begin{equation*}
V_{\text{gap}} = \max\{\text{NPV}(K^*, \phi^*) - V_{\text{AIP}},\; 0\},
\end{equation*}

where $\text{NPV}(K^*, \phi^*)$ is the net (of the full investment cost $I(K^*)$) present value of the optimally sized greenfield project. This is a comparative-statics measure of distance to optimal scale---how far the firm's gross installed-asset value falls short of the net value of a full-scale deployment---not the NPV of incremental expansion from the installed base (which would credit the firm for capacity already sunk), and not a forward-looking option value. Because the benchmark nets out the entire investment cost while assets-in-place are gross of sunk costs, the gap reaches zero before $K_{\text{installed}}$ reaches $K^*$.



\subsection*{Credit Risk Analysis}

\subsubsection*{Credit spreads}

For a levered firm with leverage $\ell$ and training fraction $\phi$, the credit spread is:

\begin{equation*}
\text{spread} = \frac{\text{coupon}}{D(X)} - r,
\end{equation*}

where $D(X)$ is the market value of debt (which differs from face value due to default risk) and the reference rate is the model's discount rate $r$ (the WACC), not a risk-free rate. Because leverage and the coupon rate are exogenous stress-test parameters (\texttt{sec-model}), the leverage gradients below describe the model's response to an imposed debt burden, not an optimally chosen capital structure. Higher leverage increases the default boundary $X_D$ and the coupon claim, reducing debt value and widening spreads. The training fraction $\phi$ enters through two opposing channels. Through the faith-based survival mechanism (Proposition 2), higher $\phi$ raises $A_{\text{eff}}$ and lowers $X_D$, reducing the \emph{probability} of default. Through the recovery specification (\texttt{sec-duopoly}), higher $\phi$ shrinks the inference business that creditors can liquidate, raising the \emph{loss given default}---the faith-based continuation value that keeps the firm alive is destroyed in bankruptcy and is therefore worthless to creditors. At the baseline calibration the loss-given-default channel modestly dominates, so spreads are mildly increasing in $\phi$ for a given leverage level, even as default probabilities fall. This two-sided interaction between training allocation and credit risk---training as solvency protection for shareholders but collateral erosion for creditors---is, to my knowledge, novel.



\subsubsection*{Default probability}

The probability of default within horizon $T$ is the first-passage probability for GBM hitting the barrier $X_D$:

\begin{equation*}
P(\text{default within } T) = \Phi(-d_1) + \left(\frac{X_D}{X}\right)^{2\nu/\sigma^2} \Phi(-d_2),
\end{equation*}

where $\nu = \mu_L - \sigma^2/2$, $d_1 = [\ln(X/X_D) + \nu T]/(\sigma\sqrt{T})$, $d_2 = [\ln(X/X_D) - \nu T]/(\sigma\sqrt{T})$, and $\Phi$ is the standard normal CDF. This is the exact barrier hitting probability for GBM, consistent with the Leland-style first-passage default mechanism in the model; the L-regime drift $\mu_L$ is the relevant one because the default boundary is an L-regime object (the firm defaults pre-switch). The calculation abstracts from the regime switch itself---a switch to $H$ before hitting $X_D$ would effectively remove default risk---so the reported figures are upper bounds on the pre-switch default hazard. Since $\mu_L$ is the risk-adjusted (certainty-equivalent) drift, the result is a \emph{model-implied risk-adjusted} first-passage probability, computed under the same convention as the spread calculation; it is not a risk-neutral (pricing-measure) probability, because the model discounts at the WACC rather than at a risk-free rate (see the footnote in \texttt{sec-model}), and the physical-measure probability (using the unadjusted drift) would be lower.



\subsection*{Dario's Dilemma}

\subsubsection*{Setup}

\subsubsection*{Value loss}

The value loss from mismatch is:

\begin{equation*}
\Delta V = \text{NPV}(\lambda_{\text{true}}, \lambda_{\text{true}}) - \text{NPV}(\lambda_{\text{true}}, \lambda_{\text{invest}}),
\end{equation*}

where $\text{NPV}(\lambda_a, \lambda_b)$ denotes the expected present value, evaluated at a common initial demand level $X_0$, of following the investment policy $(X^*, K^*, \phi^*)$ optimal for belief $\lambda_b$ when the true demand process is governed by $\lambda_a$. The timing discount $(X_0/X^*)^{\beta_H}$ adjusts for different waiting times before the respective triggers are reached. I evaluate investment \emph{outcomes} under the true demand process ($\lambda_a$) but use the \emph{investment policy} derived from the (possibly mistaken) belief $\lambda_b$. The training fraction enters the value loss through the revenue channel: a mismatched $\phi$ produces suboptimal revenue in both regimes. Throughout, percentage value losses are reported as $\Delta V / \text{NPV}(\lambda_{\text{true}}, \lambda_{\text{true}})$, the loss as a fraction of the first-best value. The levered version of the exercise keeps the same operating policy and changes only the claim being valued: the firm still chooses $(X^*, K^*, \phi^*)$ from the unlevered problem given its belief---leverage is exogenous in this model and is layered on top of an investment policy chosen on NPV grounds---and $\text{NPV}$ then denotes the value of the total claims $E + D - \ell I(K^*)$ under that policy. Because the policy is not re-optimized against the levered objective, the levered value function is not maximized at $\lambda_{\text{invest}} = \lambda_{\text{true}}$ by construction, only very nearly so: at $\ell = 0.40$ its maximizer is $\lambda_{\text{invest}} \approx 0.101$ and the value shortfall at $\lambda_{\text{true}}$ is about one thousandth of one percent, more than three orders of magnitude below the mismatch losses reported here (Internet Appendix A).



\section*{Internet Appendix}

\subsection*{A. Proofs}

\subsubsection*{Proof of Proposition 1}

*Step 1: Optimal trigger given $(K, \phi)$.* The value-matching and smooth-pasting conditions yield:

\begin{equation*}
X^*(K, \phi) = \frac{\beta_H}{\beta_H - 1} \cdot \frac{\delta K/r + cK^\gamma}{A_{\text{eff}}(\phi, K)},
\end{equation*}

where $A_{\text{eff}}(\phi, K) = [(1-\phi)K]^\alpha / (r - \mu_L + \lambda) + \lambda / (r - \mu_L + \lambda) \cdot (\phi K)^\alpha \cdot A_H$ is the effective revenue coefficient.



\emph{Step 2: NPV at the trigger.} Define $b(K) = cK^\gamma + \delta K / r$ (total cost). The revenue at the trigger is $A_{\text{eff}} \cdot X^* = \frac{\beta_H}{\beta_H - 1} b(K)$, so the NPV at the trigger is:

\begin{equation*}
V - I = \frac{1}{\beta_H - 1}\, b(K).
\end{equation*}



*Step 3: Joint optimization over $(K, \phi)$.* Substituting the trigger and NPV into $F(X_0)$:

\begin{equation*}
F(X_0) = \mathcal{C} \cdot \frac{A_{\text{eff}}(\phi, K)^{\beta_H}}{b(K)^{\beta_H - 1}} \equiv \mathcal{C} \cdot h(K, \phi),
\end{equation*}

where $\mathcal{C} = \frac{1}{\beta_H - 1}\left(\frac{(\beta_H - 1) X_0}{\beta_H}\right)^{\beta_H} > 0$ is independent of $(K, \phi)$. The optimal $(K^*, \phi^*)$ maximizes $h(K, \phi) = A_{\text{eff}}^{\beta_H} / b^{\beta_H - 1}$.



*Step 4: First-order condition for $K$.* Setting $(\partial \ln h / \partial K) = 0$:

\begin{equation*}
\tag*{[\texttt{eq-foc-K}]}
\frac{\beta_H}{A_{\text{eff}}} \frac{\partial A_{\text{eff}}}{\partial K} = \frac{(\beta_H - 1)(c\gamma K^{\gamma - 1} + \delta / r)}{cK^\gamma + \delta K / r},
\end{equation*}

Since $A_{\text{eff}}$ takes the form $g(\phi) \cdot K^\alpha$ (both in the pure H-regime and under regime switching), the left-hand side reduces to $\alpha \beta_H / K$, which is independent of $\phi$. Solving for $K$ yields the closed-form expression in Proposition 1, confirming that $K^*$ depends only on the cost and technology parameters $(\alpha, \beta_H, \gamma, \delta, r, c)$, not on the training fraction.



The first derivative with respect to $\phi$ is:

\begin{equation*}
\frac{\partial A_{\text{eff}}}{\partial \phi} = \alpha K^\alpha \bigl[-w_L (1-\phi)^{\alpha - 1} + w_H \phi^{\alpha - 1}\bigr].
\end{equation*}



\emph{Uniqueness.} The second derivative is:

\begin{equation*}
\frac{\partial^2 A_{\text{eff}}}{\partial \phi^2} = \alpha(\alpha - 1) K^\alpha \bigl[w_L (1-\phi)^{\alpha - 2} + w_H \phi^{\alpha - 2}\bigr] < 0
\end{equation*}

for all $\phi \in (0,1)$, since $\alpha \in (0,1)$ implies $\alpha(\alpha-1) < 0$ while both bracketed terms $(1-\phi)^{\alpha - 2} > 0$ and $\phi^{\alpha - 2} > 0$ are positive. Strict concavity of $A_{\text{eff}}$ in $\phi$ guarantees uniqueness.



\emph{Interior solution.} Setting $\partial A_{\text{eff}} / \partial \phi = 0$ yields:

\begin{equation*}
\frac{w_H}{w_L} = \left(\frac{1-\phi^*}{\phi^*}\right)^{\alpha - 1},
\end{equation*}

which uniquely determines $\phi^*$ as an increasing function of $w_H / w_L$. $\square$



\emph{The convention.} I nonetheless impose $A_1 = 0$, and price the pre-investment option with the pure power form $F_L(X) \propto X^{\beta_H}$ on the whole continuation region $(0, X^*)$. This is the unconditional-$A_{\text{eff}}$ convention of \texttt{sec-conventions}: option-like L-regime payoffs are discounted with the H-regime exponent $\beta_H$ rather than with a regime-switching first-passage factor mixing $X^{\beta_L^+}$ and $X^{\beta_H}$, and the exercise payoff is the unconditional L-regime installed value. With a single power in the continuation region, dividing value-matching by smooth-pasting at the investment trigger $X^*$ eliminates the coefficient and yields:

\begin{equation*}
X^* = \frac{\beta_H}{\beta_H - 1} \cdot \frac{\delta K^*/r + I(K^*)}{A_{\text{eff}}(\phi^*, K^*)},
\end{equation*}

confirming \texttt{eq-trigger-phi.} The trigger depends only on the value-matching/smooth-pasting ratio, not on the level of the coefficient, which value-matching at $X^*$ then determines; the two boundary conditions therefore pin down the trigger and the level jointly, and the system is exactly determined. The price of the convention is that this smooth-fit level differs from the forced-ODE value $C$ of \texttt{eq-particular-C}: at the baseline calibration the smooth-fit coefficient is $[V(X^*, K^*, \phi^*) - I(K^*)]/(X^*)^{\beta_H} \approx 16.5$ against $C \approx 21.1$. The direction of the resulting error is not implied by that comparison, and I therefore solve the piecewise free-boundary problem numerically and report the bias in Internet Appendix B. At the baseline calibration the exact continuation value exceeds the convention's smooth-fit value at every demand level while remaining below $C X^{\beta_H}$, so the convention understates the pre-investment option value---by 19\\% at $X_0 = X^*/2$---but the ordering is not general: it reverses for $\lambda \leq 0.02$. Propositions 1--3 are exact statements about the model these conventions define. The verification exercise shows what transfers to the exact problem and what does not: the levels of $X^*$ and $K^*$ are properties of the reduced model rather than estimates of the exact optimum, whereas the training fraction $\phi^*$ is numerically identical in the two problems and the value delivered by the reduced-form policy is within about 3\\% of the exact optimum at baseline. $\square$



\subsubsection*{Proof of Proposition 2}

Second, replacing $E_H$ with its perpetuity component:

\begin{equation*}
E_H(X) \approx A_H (\phi_i K_i)^\alpha \cdot s_i^H \cdot X - c_D/r - \delta K_i/r,
\end{equation*}

the L-regime equity ODE yields a particular solution with coefficient $A_{\text{eff},i}$, and Leland's smooth-pasting gives the default boundary below in closed form. The omitted term $D_H X^{\beta_H^-}$ is the (positive) value of the H-regime default option, so the perpetuity approximation understates the H-regime continuation value and the closed-form boundary weakly \emph{overstates} $X_D$. The magnitude of the bias is computed by solving the coupled system exactly: substituting the closed-form $E_H$ (including $D_H X^{\beta_H^-}$) into the L-regime ODE yields an additional particular term $C_2 X^{\beta_H^-}$ with $C_2 = -\lambda D_H / Q_L(\beta_H^-)$, and value-matching plus smooth-pasting at $X_D$ then determine the boundary jointly with the homogeneous coefficient. At the baseline calibration the closed-form boundary exceeds the coupled-system boundary by approximately 3\\%, uniformly across leverage levels (Internet Appendix B). The approximation therefore errs against the faith-based survival mechanism---it makes the levered firm look slightly riskier than it is---and is conservative for the credit-risk conclusions.



In the L-regime, the going-concern equity $\hat{E}(X) \equiv E(X) + (1-\ell)I(K)$ satisfies an ODE analogous to \texttt{eq-hjb-L}, with the firm defaulting when $\hat{E}(X_D) = 0$ and $\hat{E}'(X_D) = 0$ (smooth pasting). The resulting default boundary is:

\begin{equation*}
X_D = \frac{\beta_s^-}{\beta_s^- - 1} \cdot \frac{c_D/r + \delta K_i/r}{A_{\text{eff},i}},
\end{equation*}

where $\beta_s^-$ is the negative characteristic root, $c_D = c_d \cdot \ell \cdot I(K)$ is the coupon, and all comparative statics below hold *with $(K_i, \phi_i)$ fixed at their optimal values*.



*(ii) Default boundary decreasing in $\lambda$ (faith-based survival).* Holding $(K_i, \phi_i, K_j, \phi_j)$ fixed, define the L-regime revenue flow $a \equiv [(1-\phi_i)K_i]^\alpha \cdot s_i^L$ and the H-regime revenue present value $b \equiv (\phi_i K_i)^\alpha \cdot s_i^H / (r - \mu_H)$, both positive. The effective revenue coefficient is:

\begin{equation*}
A_{\text{eff},i} = \frac{a + \lambda \cdot b}{r - \mu_L + \lambda}.
\end{equation*}



Differentiating with respect to $\lambda$:

\begin{equation*}
\frac{\partial A_{\text{eff},i}}{\partial \lambda} = \frac{b(r - \mu_L) - a}{(r - \mu_L + \lambda)^2}.
\end{equation*}



This is positive whenever $b(r - \mu_L) > a$, i.e., when the H-regime revenue present value scaled by the L-regime effective discount rate exceeds the L-regime revenue flow:

\begin{equation*}
\tag*{[\texttt{eq-faith-condition}]}
\frac{(\phi_i K_i)^\alpha \cdot s_i^H}{r - \mu_H} \cdot (r - \mu_L) > [(1-\phi_i)K_i]^\alpha \cdot s_i^L.
\end{equation*}

Since $\mu_H > \mu_L$ implies $(r - \mu_L)/(r - \mu_H) > 1$, the factor $(r - \mu_L)$ amplifies the H-regime term, making the condition easier to satisfy. Intuitively, the condition requires that the training fraction $\phi_i$ is large enough that the H-regime continuation value contributes meaningfully to $A_{\text{eff},i}$.



\emph{Closed-form threshold (symmetric duopoly).} In the symmetric case with $s_i^L = s_i^H = 1/2$ and equal capacity, the contest shares and $K^\alpha$ terms cancel, and the condition reduces to $(r - \mu_L)/(r - \mu_H) \cdot \phi^\alpha > (1-\phi)^\alpha$, i.e., $(\phi/(1-\phi))^\alpha > (r-\mu_H)/(r-\mu_L)$. Solving for $\phi$:

\begin{equation*}
\phi > \underline{\phi} = \frac{R}{1 + R}, \quad R = \left(\frac{r - \mu_H}{r - \mu_L}\right)^{1/\alpha}.
\end{equation*}



For the baseline calibration ($r = 0.12$, $\mu_H = 0.06$, $\mu_L = 0.01$, $\alpha = 0.40$): $R = (0.06/0.11)^{2.5} \approx 0.22$, giving $\underline{\phi} \approx 0.18$. To complete the proof, note that $X_D$ also depends on $\lambda$ through the markup factor $M(\beta_s^-) \equiv \beta_s^- / (\beta_s^- - 1)$, since $\beta_s^-$ is the negative root of $\tfrac{1}{2}\sigma^2 \beta(\beta - 1) + \mu_L \beta - (r + \lambda) = 0$ and thus depends on $\lambda$. The full derivative decomposes into two opposing channels:

\begin{equation*}
\frac{\partial X_D}{\partial \lambda} = \underbrace{\frac{\partial M}{\partial \lambda} \cdot \frac{N}{A_{\text{eff},i}}}_{>0\;\text{($\beta$-channel)}} + \underbrace{M(\beta_s^-) \cdot N \cdot \frac{\partial (1/A_{\text{eff},i})}{\partial \lambda}}_{<0\;\text{($A_{\text{eff}}$-channel)}},
\end{equation*}

where $N = c_D/r + \delta K_i/r > 0$. The $\beta$-channel is positive because increasing $\lambda$ raises the effective discount rate $(r + \lambda)$, making $\beta_s^-$ more negative and pushing the markup $M$ toward 1 from below (since $\beta_s^- < 0$ implies $M \in (0,1)$) ($dM/d\lambda = 1/[(\beta_s^- - 1)^2 D] > 0$, where $D = \sqrt{(\mu_L - \sigma^2/2)^2 + 2\sigma^2(r+\lambda)}$). The $A_{\text{eff}}$-channel is negative under \texttt{eq-faith-condition.}



\emph{Exact net threshold.} The net sign resolves in closed form. In log form, $\partial \ln X_D / \partial \lambda = m - \partial \ln A_{\text{eff},i} / \partial \lambda$, where the markup semi-elasticity is $m \equiv (\mathrm{d}M/\mathrm{d}\lambda)/M = 1/[\beta_s^-(\beta_s^- - 1)D] > 0$, and, from the derivative above,

\begin{equation*}
\frac{\partial \ln A_{\text{eff},i}}{\partial \lambda} = \frac{b(r - \mu_L) - a}{(a + \lambda b)\,\Delta}, \qquad \Delta \equiv r - \mu_L + \lambda.
\end{equation*}



The condition $\partial X_D / \partial \lambda < 0$ is therefore $b(r - \mu_L) - a > m(a + \lambda b)\Delta$, which is \emph{linear} in $(a, b)$ and rearranges to

\begin{equation*}
\frac{b}{a} > \frac{1 + m\Delta}{(r - \mu_L) - m\lambda\Delta},
\end{equation*}

provided the regularity condition $(r - \mu_L) > m\lambda\Delta$ holds (it does throughout the calibration ranges; if it fails, no allocation can deliver $\partial X_D/\partial\lambda < 0$). With symmetric shares, $b/a = (\phi/(1-\phi))^\alpha / (r - \mu_H)$, giving the exact threshold $\tilde{\phi}$ of \texttt{eq-phi-tilde.} Since $m > 0$, $\tilde{R} > R$ and $\tilde{\phi} > \underline{\phi}$; as $m \to 0$ (no markup channel), $\tilde{\phi} \to \underline{\phi}$. At the baseline calibration $m \approx 0.77$ and $\tilde{\phi} \approx 0.32$, between $\underline{\phi} \approx 0.18$ and $\phi^* \approx 0.70$; the numerically measured $\beta$-channel is about one-third of the $A_{\text{eff}}$-channel at $\phi^*$, consistent with the closed form. The sign flip of $\partial X_D/\partial \lambda$ at $\tilde{\phi}$ is verified by finite differences as a regression test.



\subsubsection*{Proof of Proposition 3}

At zero leverage the same logic delivers an exact proof in the present model. On $(0, X_F)$ the gap admits the closed form

\begin{equation*}
L(X) - F(X) = A_{\text{eff}}^{\text{mono}}\,X - \left[\frac{\delta K_L}{r} + I(K_L)\right] - E\,X^{\beta_H}, \qquad E = \frac{V^{\text{mono}}(X_F) - V^{\text{duo}}(X_F)}{X_F^{\beta_H}} + B_F > 0,
\end{equation*}

where the first two terms are the leader's unlevered NPV of entering, the $E X^{\beta_H}$ term collects the (positive) follower-entry dilution and the follower's option coefficient, and $E > 0$ because the monopoly value exceeds the duopoly value at follower entry and $B_F > 0$. The gap is strictly concave (its second derivative is $-E\beta_H(\beta_H - 1)X^{\beta_H - 2} < 0$) and negative at $X = 0$. A strictly concave function that starts negative crosses zero from below at most once, so the up-crossing on $(0, X_L^{\text{mono}})$---which exists by the boundary conditions above---is unique. $\square$



For the follower, the cancellation is exact, by the following factorization. Write the Tullock contest payoff over the regime-relevant capacity measure $u$ (with rival measure $\bar{u}$ fixed) as $f(u) = u^{2\alpha} / (u^\alpha + \bar{u}^\alpha)$. Differentiating,

\begin{equation*}
f'(u) = \alpha\, u^{\alpha - 1}\, s\,(2 - s), \qquad s = \frac{u^\alpha}{u^\alpha + \bar{u}^\alpha}:
\end{equation*}



the marginal contest revenue is the standalone marginal revenue $\alpha u^{\alpha-1}$ scaled by the common multiplier $s(2-s)$. When $\phi_F = \phi_L = \phi$, the regime-relevant capacity ratios coincide across regimes ($u/\bar{u} = K_F/K_L$ for inference and training alike), so the share $s$---and hence the multiplier $s(2-s)$---is the \emph{same} in the L-regime and H-regime terms of the follower's allocation FOC:

\begin{equation*}
\frac{\partial A_{\text{eff},F}}{\partial \phi_F}\bigg|_{\phi_F = \phi_L = \phi} = \alpha K_F^\alpha\, s\,(2 - s)\,\bigl[-w_L (1-\phi)^{\alpha-1} + w_H \phi^{\alpha-1}\bigr].
\end{equation*}

The multiplier cancels, leaving exactly the single-firm condition of Proposition 1, Step 5, whose unique zero is $\phi^*$. Hence $\phi^*$ is an exact critical point of the follower's allocation problem for \emph{any} capacity pair $(K_F, K_L)$, not merely to first order. (The same factorization underlies the scalar reduction of the follower's capacity problem in Internet Appendix B.) What remains computational is global optimality: that this critical point is the follower's global maximizer is confirmed numerically---$\phi_F^* = \phi_L^* = \phi^*$ holds to at least four decimal places across the full parameter space used in the calibration (varying $\sigma$, $\mu_H$, $\alpha$, $\gamma$, $\lambda$, and $\ell$ over their calibration ranges), with no parameterization showing an economically meaningful gap. The duration of the monopoly phase is substantial---$X_F^*/X_P \approx 44$--$46$ at baseline, roughly invariant to leverage---but it operates on the \emph{timing} margin (part iii), not the allocation margin. $\square$



\subsubsection*{Numerical Finding 1 (Dario's dilemma)}

\subsection*{B. Numerical Verification Methods}

\textbf{Coupled default boundary.} The single-boundary default formula (\texttt{eq-default-boundary}) is verified against the fully coupled regime-switching system: the H-regime equity has the closed Leland form $E_H(X) = A_H(\phi K)^\alpha s^H X - (c_D + \delta K)/r + D_H X^{\beta_H^-}$, whose default-option term forces an additional particular solution $C_2 X^{\beta_H^-}$ in the L-regime equity ODE; value-matching and smooth-pasting at the boundary then determine $X_D$ jointly with the homogeneous coefficient $A_-$. Because $A_-$ enters both boundary conditions linearly, the combination $\beta_L^- E(X_D) - X_D E'(X_D) = 0$ eliminates it exactly, reducing the coupled system to a single equation,

\begin{equation*}
A_{\text{eff}}\,(\beta_L^- - 1)\,X_D - \beta_L^- N + C_2\,(\beta_L^- - \beta_H^-)\,X_D^{\beta_H^-} = 0, \qquad N = \frac{c_D + \delta K}{r},
\end{equation*}



whose economically relevant root (the continuation of the single-boundary root as $C_2 \to 0$) is computed by Brent's method. A first-order expansion in the small coefficient $C_2$ around the single-boundary root $X_D^0$ gives the closed-form relative bias

\begin{equation*}
\frac{X_D^0 - X_D}{X_D^0} \;\approx\; \kappa = \frac{C_2\,(\beta_L^- - \beta_H^-)\,(X_D^0)^{\beta_H^- - 1}}{A_{\text{eff}}\,(\beta_L^- - 1)} > 0,
\end{equation*}

positive because $C_2 < 0$ at calibration values---so the direction of the bias (the single-boundary formula overstates default risk) follows analytically, not only numerically. At baseline the exact relative bias is 3.1\\% and the first-order $\kappa$ is 2.9\\%, uniform across leverage levels; the closed-form boundary exceeds the coupled-system boundary by approximately 3\\% throughout, in the conservative direction.



\textbf{Piecewise low-regime stopping problem.} The pure-power convention for the pre-investment option (Proof of Proposition 1, Step 5b) is verified against the exact piecewise free-boundary problem, in the same spirit as the coupled default boundary above. Below $X_H^*$ the forcing term in \texttt{eq-hjb-L} is the H-regime option value and the continuation value is $A_1 X^{\beta_L^+} + C X^{\beta_H}$, the negative root being removed by $F_L(0) = 0$. Above $X_H^*$ a switch is followed by immediate exercise, so the forcing is affine, $\lambda\,[a_H X - b_H]$ with $a_H = A_H (K_H^*)^\alpha$ and $b_H = \delta K_H^*/r + I(K_H^*)$, and the continuation value is

\begin{equation*}
D_1 X^{\beta_L^+} + D_2 X^{\beta_L^-} + \frac{\lambda a_H}{r + \lambda - \mu_L}\,X - \frac{\lambda b_H}{r + \lambda}.
\end{equation*}

Value and slope matching at $X_H^*$, together with value matching and smooth pasting at the free boundary, determine $(A_1, D_1, D_2, X^*)$: the coefficients solve a linear system given $X^*$, and the remaining slope condition at $X_H^*$ is solved by Brent's method. When the resulting boundary falls below $X_H^*$ the problem collapses to the single-region case, in which the pure-power forcing is exact. The solution is implemented in the replication package (module `piecewise\_option.py`) and cross-checked against an independent Brennan--Schwartz solution of the associated linear complementarity problem on a log-demand grid, which reproduces the free boundary to a relative $10^{-4}$ and the value function to a relative $10^{-9}$.



\textbf{Separable reduction of the follower problem.} At a common training fraction $\phi_F = \phi_L = \phi$, the Tullock denominators factor---$[(1-\phi)K_F]^\alpha + [(1-\phi)K_L]^\alpha = (1-\phi)^\alpha (K_F^\alpha + K_L^\alpha)$, and likewise for the training measure---so the follower's effective revenue coefficient separates exactly as in Proposition 1:

\begin{equation*}
A_{\text{eff},F} = g(\phi) \cdot \frac{K_F^{2\alpha}}{K_F^\alpha + K_L^\alpha},
\end{equation*}



with $g(\phi) = w_L (1-\phi)^\alpha + w_H \phi^\alpha$. The allocation FOC is then exactly the single-firm condition (the role-invariance of Proposition 3(ii)), and the capacity FOC reduces to a single equation with \emph{effective revenue elasticity} $\alpha(2 - s_F)$, where $s_F = K_F^\alpha/(K_F^\alpha + K_L^\alpha)$:

\begin{equation*}
\alpha \beta_H\,\bigl(2 - s_F(K_F)\bigr)\, b(K_F) = (\beta_H - 1)\, K_F\, b'(K_F),
\end{equation*}

where $b(K) = \delta K / r + c[1 - \ell(1 - c_d/r)]K^\gamma$ is the leverage-adjusted total cost. The left-hand bracket $\alpha\beta_H(2 - s_F) - (\beta_H - 1)Kb'/b$ is strictly decreasing in $K$ ($s_F$ is increasing and $Kb'/b$ rises from 1 to $\gamma$), so the interior optimum is \emph{unique} whenever $(\beta_H - 1)/(2\alpha\beta_H) < 1$ and $(\beta_H - 1)/(\alpha\beta_H) > 1/\gamma$. The first condition is weaker than (A2)'s upper bound because the follower's small-$K$ elasticity is $2\alpha$ (share gains compound the standalone revenue gain), so the follower's problem remains well-posed in part of the region where the single-firm problem degenerates. The scalar solution agrees with the two-dimensional Nelder-Mead optimizer to a relative tolerance of $10^{-6}$ at baseline (with and without leverage), which is enforced as a regression test; the elasticity wedge $\alpha(2 - s_F) > \alpha$ is the analytical source of the large follower scale ($K_F \gg K_L$) noted in \texttt{sec-calibration.}



\end{document}
